\section{An MVCC-Consistent CSR Index}
\label{sec:index}

\todo{2.5 pages. Figure: epoch + change log + per-reader snapshot.}
\subsection{Design constraints from the extension interface}
\todo{What the index type interface delivers (appends, deletes with row ids, serialize, WAL replay) and what it does
not (commit ids, visibility); keys not row ids identify vertices; the index may not evaluate visibility at past
bounds because the engine garbage-collects row versions.}
\subsection{Epochs, change groups and snapshots}
\todo{Immutable epoch CSR with dictionary and row ids per position; change log stamped per commit group; a reader
decides visibility once per group with CanFetch; snapshot = epoch partitions + delta partitions + suppressed bitmap;
private snapshot for writing transactions; snapshot cache keyed by (generation, newest visible group).}
\subsection{Fold}
\todo{Rebuilding an epoch from a live transaction, generation stamping, retained epochs, id re-resolution; fold
ratio policy.}
\subsection{Checkpoint, recovery and parallel load}
\todo{Segments in the index allocator, manifest in storage options, WAL replay through Append/Delete, parallel
deserialization numbers (SF300 1.15 s to 0.09 s).}
\subsection{Path reconstruction}
\todo{Parent pointers in the single-source traversal, edge row ids from epoch and delta partitions, keys for vertices;
what the multi-source traversal cannot do and why.}
\subsection{Correctness}
\todo{Argument: a reader sees exactly the edges whose rows are visible to its transaction; cases: concurrent insert,
delete of an epoch edge, delete of a delta edge, rollback, fold during a long snapshot, restart.}
